\chapter{Limites de l'analyse}
\label{chap:limites}

\begin{constraints}
\begin{itemize}
\item Les données sont observationnelles.
\item L'analyse décrit des associations et des distributions ; elle ne permet pas d'établir une causalité.
\item Le jeu est très équilibré : les volumes par modalité sont proches, ce qui limite la visibilité d'éventuels contrastes.
\item Les différences de KPI entre catégories sont souvent très faibles ; une échelle trop zoomée les exagérerait.
\item Les corrélations CTR / clics / impressions sont en partie mécaniques, car le CTR est un ratio de ces deux grandeurs.
\item Les résultats sont limités au dataset étudié (200\,005 campagnes de 2021) et ne doivent pas être généralisés au-delà.
\item Le CTR moyen des campagnes (moyenne des ratios) n'est pas égal au CTR global ($\sum$ clics / $\sum$ impressions) ; les deux indicateurs répondent à des questions différentes.
\end{itemize}
\end{constraints}

Aucune recommandation d'allocation budgétaire, de ciblage ou de choix de canal n'est formulée à partir de ces écarts.
